% ============================================================
% 第3章：关键算法与核心原理
% ============================================================

\section{四级容错需求解析链}

需求解析模块将自然语言选型需求转换为 \texttt{RequirementConstraints} 结构化约束对象。针对工程需求文本表述多样、参数边界隐含的特点，系统设计了四级串联容错策略：前级完成基础语义映射，后级针对已知缺陷进行局部修正，在保留 LLM 语义理解能力的同时增强规则约束下的工程可靠性。具体机制如下：

\textbf{第一级——LLM语义解析}：调用 LLM 从输入文本提取 \texttt{category}、\texttt{topology}、输入/输出电压与电流、温度范围及等级偏好等字段，以 JSON 格式返回。

\textbf{第二级——规则化覆盖}：通过关键词匹配修正 LLM 输出——"降压""buck""DC-DC"触发 \texttt{dc\_dc\_converter}，"升压""boost"触发对应拓扑；同时对否定句式（如"非车规"）进行专项处理，正确映射为工业级。

\textbf{第三级——电压方向纠正}：DC-DC 类别且输入输出电压均已提取时，根据两者大小关系强制修正拓扑——$V_{\text{in}} > V_{\text{out}}$ 设为 \texttt{buck}，反之为 \texttt{boost}。

\textbf{第四级——温度范围正则匹配}：通过多格式正则（"-40到85°C""-40 to 125°C""-40\textasciitilde{}85"）提取温度边界；电流提取优先匹配 mA 单位并除以 1000，避免"500 mA"被误识别为 500 A 的量级错误。


\section{eZ-PLM API 集成与多前缀分组查询}

\subsection{HMAC-SHA256 签名认证}

eZ-PLM API 采用 HMAC-SHA256 签名认证。每次请求需携带 API 密钥、Unix 时间戳、UUID 随机数和签名值。签名内容由请求方法、路径、规范化查询参数、时间戳和随机数按固定顺序拼接形成待签名字符串，以 API 密钥为 HMAC 种子计算 HMAC-SHA256 并进行 Base64URLSafe 编码。客户端每次请求生成新的时间戳和随机数，以降低重放攻击风险。

\subsection{多前缀分组查询策略}

由于 eZ-PLM API 主要支持基于制造商型号前缀的关键字检索，系统构建了从类别/拓扑到型号前缀的映射机制（表\ref{tab:keywords}）：Buck 映射至 TPS54、LM2596 等前缀，Boost 映射至 TPS61、MCP1640 等，LDO 映射至 TPS79、MCP1703 等。系统对每个前缀独立发起查询（每前缀最多 50 条），以型号为键去重，总召回上限 200 条。

\begin{table}[H]
\centering
\caption{类别/拓扑至制造商型号的前缀映射}
\label{tab:keywords}
\footnotesize
\setlength{\tabcolsep}{4pt}
\renewcommand{\arraystretch}{0.9}
\begin{tabularx}{\textwidth}{p{2.8cm}p{1cm}X}
\toprule
\textbf{类别} & \textbf{拓扑} & \textbf{查询前缀列表} \\
\midrule
\texttt{dc\_dc\_converter} & \texttt{buck} & \texttt{TPS54}、\texttt{TPS62}、\texttt{LM2596}、\texttt{LM2576}、\texttt{ADP23}、\texttt{LTC388}、\texttt{LTC365}、\texttt{ST1S}、\texttt{L5970} \\
\texttt{dc\_dc\_converter} & \texttt{boost} & \texttt{TPS61}、\texttt{TPS63}、\texttt{LTC370}、\texttt{LTC358}、\texttt{MCP1640} \\
\texttt{ldo} & --- & \texttt{TPS79}、\texttt{TPS72}、\texttt{MCP1703}、\texttt{MCP1700}、\texttt{ADP312}、\texttt{LT1763}、\texttt{MCP1501} \\
\bottomrule
\end{tabularx}
\end{table}

针对 API 返回结果中输出电压字段可能混用固定值与最大电压范围的问题，系统设计了基于 MPN 的电压推断方法，通过正则从型号字符串提取固定输出电压（如 \texttt{LM2596S-5.0} $\rightarrow$ 5.0 V）并在匹配阶段优先采用。数据检索采用“API 优先、离线数据库兜底”策略，候选器件经五维过滤（类别匹配、输入电压覆盖、输出电压 $\pm 5\%$ 容差、输出电流下限、温度范围/等级匹配）后保留满足全部约束的器件。


\section{双模自适应混合评分引擎}

评分引擎的总体架构（双模切换逻辑、权重公式及推荐分级策略）已在第 2.3 节论证。本节重点介绍各子维度计算逻辑、LLM 提示词设计及数值边界处理。

\subsection{参数匹配评分的逐维计算}

参数匹配分 $S_{\text{param}}$ 基于实际约束项数量 $N$ 动态归一化，包含四个子维度：

\begin{itemize}
\item \textbf{输入电压覆盖}：器件输入电压范围包含用户标称值（$V_{\text{in,min}} \le V_{\text{in,nominal}} \le V_{\text{in,max}}$）得 1，否则 0。
\item \textbf{输出电压匹配}：固定输出器件通过 MPN 推断提取额定值，与需求偏差小于 0.01 V 得 1.0，$\le$ 5\% 得 0.6，超过 5\% 得 0.1；ADJ 器件不计入 $N$。
\item \textbf{输出电流满足}：$s_{\text{iout}} = \min(1.0, I_{\text{out,part}} / (2 \times I_{\text{out,demand}}))$，器件电流达需求两倍时满分，恰好满足时得 0.5。
\item \textbf{温度范围覆盖}：器件温度区间完全覆盖需求区间得 1，否则 0。未指定温度约束则该维度不计入评分。
\end{itemize}

\subsection{供应链稳定性评分}

供应链稳定性分 $S_{\text{supply}}$ 基于库存与生命周期状态分段计算：库存 $>1000$ 件得 100 分，100–1000 得 70，1–100 得 40，0 件得 10；若生命周期为 \texttt{obsolete} 或 \texttt{discontinued}，库存得分上限截断为 10 分。

\subsection{LLM 软评分提示词与安全边界}

LLM 增强模式下，\texttt{score\_part\_with\_llm()} 构造系统提示词（角色定位为“资深电子工程师”）与用户提示词（含原始需求、器件电气参数及最多三条参考设计案例摘要），要求返回 JSON 格式评分 \texttt{\{"application\_score": 85, "design\_risk\_score": 78\}}。系统对返回执行双阶段安全处理：先截取首尾大括号内 JSON 解析，再对评分值执行 $\max(0, \min(100, v))$ clamp；解析失败或字段缺失则回退至纯规则模式。

\section{双层风险评估架构}

风险报告采用规则验证与 LLM 叙述分离的双层设计。\textbf{第一阶段——规则引擎检测}：规则引擎对供应链与工程两类风险进行确定性检测。供应链风险包括无候选器件、推荐不足、供应商过度集中、库存异常（零库存或低于 100 件）、停产状态等；工程风险包括车规未满足、参数匹配度低、关键字段缺失等。每条检测结果按高/中/低三级标注严重程度并附缓解建议，形成结构化 \texttt{RiskItem} 列表。

\textbf{第二阶段——LLM 背景说明}：已验证的 \texttt{RiskItem} 列表传入 LLM，仅允许对已确认风险提供工程背景说明，不可独立产生新风险结论，且每条文本显式标注数据来源为“规则引擎检测结果”，实现风险事实与语言表述的完全分离。报告遵循 ISO 31000\cite{iso31000} 框架，包含执行摘要、$5\times 5$ 风险矩阵\cite{iec60812}、风险详述、人工复核清单及结论建议。

\section{检索增强生成（RAG）工程知识库}

RAG 模块基于 ChromaDB 持久化向量存储及 \texttt{all-MiniLM-L6-v2} 嵌入模型（384 维）构建\cite{lewis2020rag}，当前包含 29 条工程知识条目，覆盖 Buck/Boost/LDO 设计公式、热管理、车规 AEC-Q100 认证\cite{aecq100}、PCB Layout 规范\cite{ipc2221,ipc2152}、EMI/可靠性指南等 11 个类别。每次选型分析时自动检索 Top-5 相关条目，以 \texttt{EvidenceIR} 形式注入证据链，并在三类输出报告中展示相关知识、相关度评分及来源引用。
